\documentclass[11pt]{article}
\usepackage{amsmath,amssymb}
\usepackage[margin=1in]{geometry}

\title{Structure of the STWO Private-Witness ZK Integration}
\author{Staff Crypto Researcher \and Staff Security Researcher}
\date{}

\begin{document}
\maketitle

\section{Status and Scope}

This note describes the reviewed structure of the current STWO private-witness
zero-knowledge integration. It is an integration note, not a production-readiness
claim. The public STWO prover and verifier remain unchanged; the ZK path is
additive and is guarded by verifier-owned metadata checks.

The integration currently covers private witness randomization, verifier-bound
privacy metadata, dynamic ZK degree accounting, quotient/composition masking
metadata, FRI sample planning guardrails, and target-bearing private LogUp policy.
Private LogUp is accepted only when hidden claims have a verifier-checkable
algebraic target. AIRs without such a target remain fail-closed.

\section{Code Structure}

The integration is organized around three concepts.

\paragraph{AIR privacy provider.}
An AIR-specific provider declares which committed columns are public or private,
which private roots imply dependent private columns, which semantic trace domains
those columns must preserve, and which LogUp claims are semantically public,
target-bearing private, or unsupported. The framework does not guess privacy from
column names; the AIR declares the policy and the verifier binds the same policy.

\paragraph{Canonical ZK metadata builder.}
The generic builder canonicalizes the AIR provider output into verifier-owned
metadata: privacy-map hash, public-statement hash, private-column scope hash,
randomizer-space hash, per-column degree overlays, quotient split mask profile,
LogUp claim manifest, and statistical-budget records. This is the mechanism that
turns AIR-specific privacy declarations into a reusable STWO ZK plan.

\paragraph{Masked proving and audited verification.}
The prover uses the canonical metadata to randomize private witness polynomials
and to build masked opening data. The verifier receives or derives the expected
metadata from verifier-owned configuration, not from the proof alone, and rejects
proofs whose privacy map, degree profile, sample-point plan, quotient profile,
LogUp policy, or statistical budget do not match that configuration. Example AIR
wrappers derive this configuration from their provider state; the generic API
therefore requires callers to preserve that verifier-owned sequencing discipline.

\section{Witness Masking}

For a private witness polynomial \(w(X)\) over semantic trace domain \(H\), the
masking follows the standard vanishing-polynomial construction:

\[
    \widehat{w}(X) = w(X) + Z_H(X) r(X),
\]

where \(Z_H(X)\) vanishes on \(H\), and \(r(X)\) is fresh prover randomness for
the proof. On every point in \(H\), the masked polynomial agrees with the original
witness:

\[
    x \in H \implies \widehat{w}(x) = w(x).
\]

Away from \(H\), openings are statistically hidden by the randomizer, subject to
the declared degree bound and the field-size security budget.

Some STWO traces use private columns that participate in more than one semantic
domain. In that case the implementation uses the vanishing polynomial for the
declared domain set:

\[
    \widehat{w}(X)
    =
    w(X)
    +
    \left(\prod_{H_i \in \mathcal{H}_{\mathrm{sem}}} Z_{H_i}(X)\right) r(X).
\]

This preserves all declared semantic-domain evaluations while still randomizing
off-domain openings. The degree increase is derived from the total vanishing
degree of the declared domain set rather than from a hardcoded AIR constant.

\section{Degree Accounting}

The masked witness degree is not treated as a fixed offset. The ZK degree profile
is derived from:

\[
    \deg(w),\quad
    \deg\!\left(\prod_i Z_{H_i}\right),\quad
    \deg(r),\quad
    \text{AIR constraint expansion},\quad
    \text{composition split},\quad
    \text{FRI blowup}.
\]

The verifier binds the resulting per-column overlays and quotient split profile.
This prevents an implementation from silently shrinking the randomizer space,
omitting a private dependency, or using a proof-local quotient layout that the
verifier did not authorize.

\section{Quotient and Composition Masking}

The composition quotient is evaluated against the randomized private trace. The
quotient is then split into STWO composition columns as usual, but the ZK path
binds the split layout and expected mask profile into verifier metadata. The
important invariant is that quotient masking is not a prover-only decoration:
the verifier checks that the degree and split profile match the ZK metadata it
derived before accepting the proof.

\section{FRI Sample Planning}

STWO can open columns that live on different committed domains. The ZK path uses
a semantic sample-point plan that records how each committed-domain point relates
to the AIR semantic point and to the FRI lifting layer. For a committed point
\(p_i\), semantic point \(z_i\), committed log size \(c_i\), and FRI lifting log
size \(L\), the verifier-owned plan enforces the relation

\[
    \operatorname{double}^{L-c_i}(p_i) = z_i.
\]

This prevents private masking from using sample points that are inconsistent with
the original AIR semantics or with the FRI query plan.

\section{Target-Bearing Private LogUp}

LogUp introduces a separate privacy issue. A LogUp claimed sum can itself be a
witness-derived fingerprint. Hiding it locally is not enough: if the verifier no
longer checks an algebraic target, lookup soundness can become vacuous.

The accepted policies are:

\[
    \mathsf{SemanticallyPublic},\qquad
    \mathsf{StatisticalAggregate}\{\mathsf{aggregate\_id}\},\qquad
    \mathsf{PrivateUnsupported}.
\]

A target-bearing private aggregate must have a verifier-checkable target

\[
    T \in \{0,\ \mathsf{PublicExpression}\}.
\]

For private claims \(s_i\), masked claims \(m_i\), correction shifts \(c_i\), and
claim-domain sizes \(|H_i|\), verification checks

\[
    \sum_i \left(m_i - |H_i| c_i\right) = T.
\]

The correction hides each private claim while preserving the public aggregate
relation. If no valid target exists, the AIR remains fail-closed. This is why
standalone Poseidon private LogUp is not enabled by the current integration: its
standalone claimed sum is witness-derived and does not currently have a reviewed
zero, public-expression, or aggregate target.

\section{Transcript Binding}

The verifier binds ZK metadata before challenges that can be affected by private
data. In the core ZK path, metadata is bound before composition, OODS, sampled
value, and FRI-query challenges. AIR-specific wrappers must additionally bind the
same verifier-owned metadata before any lookup challenge they draw outside the
core verifier. The relevant challenge boundaries are commitments, lookup
challenges, composition challenges, OODS sampling, sampled values, and FRI
queries. The security property is that the prover cannot choose a privacy map,
randomizer space, LogUp policy, or sample-point layout after seeing challenges
that would make that choice favorable.

\section{What This Improves in STWO}

The integration improves STWO by adding a reusable private-witness framework
instead of one-off masking for a single example. The main improvements are:

\begin{itemize}
    \item verifier-owned privacy metadata and canonical hashes;
    \item AIR-declared public/private column policy;
    \item private witness masking over STWO semantic domains;
    \item support for mixed semantic-domain masking;
    \item dynamic degree-bound derivation from the declared domain set;
    \item quotient split mask-profile binding;
    \item FRI sample-point consistency checks;
    \item target-bearing private LogUp policy;
    \item fail-closed handling for AIRs without verifier-checkable LogUp targets;
    \item additive compatibility with ordinary public STWO proving.
\end{itemize}

\section{Current Limitations}

This is not a claim that every STWO AIR is already private. Each AIR must still
declare a privacy provider and, for private LogUp, a verifier-checkable target.

Standalone Poseidon private LogUp remains fail-closed. Exposing its raw
claimed-sum would leak a witness-derived fingerprint, while hiding it without a
target would fail to prove lookup balance. Poseidon can be enabled only after a
reviewed zero, public-expression, or aggregate target is derived for the concrete
statement being proved.

Benchmark artifacts and baseline reports are intentionally not part of this note.
Performance measurement should be generated on demand from benchmarks, not kept
as stale repository documentation.

\section{Conclusion}

The current integration maps the paper-shaped vanishing-polynomial masking idea
onto STWO's committed trace, quotient, and FRI infrastructure. The key design
choice is that privacy is verifier-owned and target-bearing: the prover may mask
private witness material, but the verifier still binds the public statement,
degree budget, sample plan, quotient profile, and LogUp algebraic target before
accepting.

\end{document}
